\begin{afterword}
\summaryhead

The essay starts from Scott Aaronson's remark that libertarianism and the many-worlds interpretation of quantum mechanics are both cases of ``bullet-swallowing.'' The author proposes a different analogy, between libertarianism and Science, by which he means an idealized scientific method rather than the practice of scientists.

Both, he says, distrust reasonable-sounding arguments, try to build systems more trustworthy than their members, and harness human flaws. Markets turn selfishness into production. Science turns each scientist's stubborn wish to prove a theory into a motive to run experiments that could refute it, with replication as a check. Both depend on people being somewhat better than selfish, and both degrade gracefully when they are not.

The author concludes that Science is not easily reconciled with probability theory. It does not trust individuals to reason correctly and demands a decisive experiment instead. Letting Bayesian reasoning override Science is ``not a trivial step.'' The essay ends by asking whether the reader will still believe in wavefunction collapse.

\responsehead

The essay's argument and its purpose point in opposite directions. The argument is a case for institutional humility: individuals reason badly, reasonable-sounding arguments mislead, and a system built on that assumption is ``somewhat more trustworthy than its parts.'' The purpose was stated the day before, in ``The Dilemma: Science or Bayes?'': to ``break your allegiance to Science'' over many-worlds. The last paragraph pursues it with a dare: ``do you think you're smarter than that?''

Nothing in the essay supports a yes. By its first principle, the author's Occam's-razor case for many-worlds is the kind of reasonable-sounding argument Science is right to distrust. By its own rule about stolen predictions, Science is right to withhold credit from a later theory that, in the author's words, ``doesn't seem to make any new predictions relative to the old theory.''

The route to the dare is built from borrowed and unsourced material. Aaronson's post is quoted for its setup and not for its point, which was that one ``would've been wise to doubt'' the bullet-swallowers' conclusions even without spotting a flaw. The account of Science as a system that harnesses self-interest is Polanyi's (1962) and Kitcher's (1990), uncredited. The history comes without dates or sources. The libertarian half of the analogy is a contested political program presented as settled history, and it is exempted from its own distrust of lovely theories.

The definitions do quiet work. ``Science'' is capitalized and made an ideal, so no scientist or institution can contradict what the author says it holds. The model omits peer review, the institution that most resembles what the essay says Science refuses to do. ``Rational'' is defined as the output of a correct probability calculation, so the claim that a correct calculation gives the rational answer is true by definition. It says nothing about whether anyone's calculation is correct.

The essay offers its reader two identities: the humble member of a system that distrusts itself, and the individual smart enough to overrule it. It spends its length on the first and closes by rewarding the second.

In short: an argument for distrusting clever individual reasoning, assembled from uncredited sources, that ends by inviting the reader to trust their own over Science on the author's say-so.
\end{afterword}
